\documentclass[11pt,a4paper]{article}

\usepackage[T1]{fontenc}
\usepackage[utf8]{inputenc}
\usepackage{lmodern}
\usepackage{microtype}
\usepackage[a4paper,margin=27mm,headheight=14pt]{geometry}
\usepackage{amsmath,amssymb,amsthm,mathtools}
\usepackage{booktabs,longtable,array,tabularx}
\usepackage{float}
\usepackage{enumitem}
\usepackage{xcolor}
\usepackage{fancyhdr}
\usepackage{hyperref}

\definecolor{deepblue}{RGB}{35,72,112}
\definecolor{softblue}{RGB}{238,245,251}
\definecolor{softgray}{RGB}{245,246,248}
\definecolor{darkgray}{RGB}{70,74,80}

\hypersetup{
  colorlinks=true,
  linkcolor=deepblue,
  citecolor=deepblue,
  urlcolor=deepblue,
  pdftitle={From Bars to Analytic HFT Support},
  pdfauthor={Wiser Massive Modes research project}
}

\pagestyle{fancy}
\fancyhf{}
\fancyhead[L]{\small From Bars to analytic HFT support}
\fancyhead[R]{\small TK6B.5f}
\fancyfoot[C]{\thepage}

\setlist[itemize]{leftmargin=1.6em,itemsep=0.25em,topsep=0.4em}
\setlist[enumerate]{leftmargin=1.8em,itemsep=0.3em,topsep=0.4em}
\setlength{\parindent}{0pt}
\setlength{\parskip}{0.55em}
\renewcommand{\arraystretch}{1.15}

\newtheorem{theorem}{Theorem}[section]
\newtheorem{lemma}[theorem]{Lemma}
\newtheorem{proposition}[theorem]{Proposition}
\newtheorem{corollary}[theorem]{Corollary}
\theoremstyle{definition}
\newtheorem{definition}[theorem]{Definition}
\newtheorem{remark}[theorem]{Remark}

\newcommand{\Spin}{\operatorname{Spin}}
\newcommand{\SO}{\operatorname{SO}}
\newcommand{\PE}{\operatorname{PE}}
\newcommand{\PEF}{\operatorname{PE}_{F}}
\newcommand{\Res}{\operatorname{Res}}
\newcommand{\Rep}{R}
\newcommand{\Kseven}{K_{7}}
\newcommand{\cF}{\mathcal{F}}
\newcommand{\one}{\mathbf{1}}

\newenvironment{summarybox}
  {\begin{center}\begin{minipage}{0.92\textwidth}\color{darkgray}
   \hrule\vspace{0.7em}}
  {\vspace{0.5em}\hrule\end{minipage}\end{center}}

\title{\textbf{From Bars to Analytic HFT Support}\\
\large The trajectory--kernel route to the TK6B.5f all-orders theorem}
\author{Wiser Massive Modes research project}
\date{1 October 2026}

\begin{document}

\maketitle

\begin{abstract}
Bars proposed that at strong coupling the perturbative massive primaries of
the type-IIA superstring should be supplemented by nonperturbative states so
that their cumulative $\Spin(9)$ content assembles into $\Spin(10)$
representations.  In its minimal form the claim is an existence problem: for
every level $n$, the cumulative character
$R_n^9=\sum_{i=1}^{n}r_i^9$ should admit at least one positive integral
$\Spin(10)$ lift.  The lift need not be unique.

This report first develops the full route from Bars's proposal to the local
all-orders theorem proved in TK6B.5f.  It reviews the covariant HFT generating
function, multiplicity-free $D_5\downarrow B_4$ branching, Regge shifts,
restriction-kernel coordinates, positive packet cones, Gale circuits,
interlacing shells, the local M\"obius inverse and stable Pieri translation.
The mathematical notions beyond a standard graduate representation-theory
course are introduced when they enter.

The second part proves the HFT support theorem.  A finite triality identity
telescopes the Green--Schwarz oscillator product and performs the
universal-supermultiplet division before support estimates are applied.  A
Weyl-numerator argument transfers $\Spin(8)$ support to $\Spin(9)$, uniform
vector and spinor estimates close the infinite oscillator tail, and a fixed
323-record low-mode certificate finishes the proof.  This closes the HFT
premise needed for stable Pieri transport at every degree.  The remaining
Bars problem is global: control all moving weighted circuits and prove
integral packetization at arbitrary level.
\end{abstract}

\begin{summarybox}
\textbf{Three levels of the argument.}

\emph{Bars's target:} for every $n$, find a positive integral $\Spin(10)$
character $X_n$ with
\[
   X_n\mathord\downarrow_{\Spin(9)}=R_n^9
   =\sum_{i=1}^{n}r_i^9.
\]
\emph{The theorem proved here:} if the irreducible $\Spin(9)$ character
$V_\mu$ occurs with nonzero signed coefficient in degree $d$ of the virtual
HFT remainder $F(q)$, then
\[
   [V_\mu]F_d\neq0 \quad\Longrightarrow\quad d\geq K(\mu),
\]
where $K$ is the exact Pieri-hazard threshold derived in TK6B.5e.  The theorem
holds at every oscillator degree.

\emph{The remaining gap:} show that stable transport controls every
horizon-dependent weighted circuit and supply an all-level integral packet
decomposition (or a saturation theorem).  Finite levels are evidence and
regression data, not the all-orders argument.
\end{summarybox}

\tableofcontents
\newpage

\section{Bars's strong-coupling proposal}

\subsection{Why $\Spin(9)$ and $\Spin(10)$ appear}

For a massive particle in ten-dimensional Minkowski space the little group
is $\Spin(9)$.  Perturbative massive string states therefore decompose into
$\Spin(9)$ irreducible representations.  If strongly coupled type-IIA string
theory develops an eleventh spacetime dimension, the corresponding massive
rotation group is $\Spin(10)$.  A necessary representation-theoretic shadow
of that claim is that the ten-dimensional states, after adding appropriate
nonperturbative states, should be restrictions of $\Spin(10)$
representations.

The supersymmetry descendants are conveniently factored off.  Let $r_n^9$
denote the $\Spin(9)$ character of the chiral massive \emph{primaries} at
perturbative level $n$.  Acting with the broken supercharges produces the
universal chiral massive multiplet
\[
 Z_Q=[2,0,0,0]_9+[1,0,0,1]_9+[0,0,1,0]_9,
\]
of dimension $44+128+84=256$.  For the closed type-II theory there is one
such factor from each chiral half.  Bars's question concerns the remaining
primary factor: can it also be organized covariantly with respect to the
extra spatial direction?  This is the content relevant to the present
project, rather than a claim that the perturbative spectrum alone already
has $\Spin(10)$ symmetry.

We use $\Spin$ rather than $SO$ when spinor representations matter.  The
older literature often writes $SO(9)$ and $SO(10)$ for the same character
problem.

\subsection{Bars's additional level and the cumulative sum}

Bars introduces an elementary string level $l$ and an additional
nonperturbative level $l'$.  At fixed ten-dimensional mass level $n$, his
on-shell proposal is
\[
   n=l'+l,
\]
with the pairs
\[
   (l',l)=(0,n),(1,n-1),\ldots,(n-1,1).
\]
The $\Spin(9)$ quantum numbers depend on $l$, while $l'$ distinguishes the
new states.  In the single-multiplicity, or ``minimal'', scheme each pair is
used once.  Consequently one chiral half at mass level $n$ has the cumulative
primary character
\begin{equation}
\label{eq:bars-cumulative}
   R_n^9=\sum_{l=1}^{n}r_l^9.
\end{equation}
Bars checked directly through $n=5$ that these cumulative characters
reassemble into $\Spin(10)$ characters \cite{Bars}.  For orientation, his
dimension notation gives the following first examples.  Bosonic and
fermionic sectors are displayed separately where needed.

\begin{table}[H]
\centering
\caption{Bars's minimal cumulative pattern for one chiral half.  These are
the examples displayed in the 1995 paper, not an all-orders proof.}
\label{tab:bars-low-levels}
\begin{tabularx}{\textwidth}{cXl}
\toprule
$n$ & cumulative $\Spin(9)$ primary content & one $\Spin(10)$ parent \\
\midrule
1 & $1$ & $1$ \\
2 & $1+9$ & $10$ \\
3 & $(1+9+44)_B+16_F$ & $54_B+16_F$ \\
4 & $(1+2\cdot9+36+156+44)_B+(16+128)_F$
  & $(45+210)_B+144_F$ \\
5 & $\bigl(2(1+9+36+44)+84+156+231+450\bigr)_B$
    $+\bigl(2(16+128)+576\bigr)_F$
  & $(1+120+320+660)_B+(144+720)_F$ \\
\bottomrule
\end{tabularx}
\end{table}

The representation dimensions in the last column suppress Dynkin labels and,
for spinors, an outer-chirality choice.  Later finite computations show that
positive parents need not be unique.  Bars's Section III.B claim therefore
has the following precise minimal interpretation.

\begin{definition}[Minimal Bars lifting problem]
For every $n\geq1$, find at least one positive integral $\Spin(10)$
character $X_n$ such that
\[
   \Res^{\Spin(10)}_{\Spin(9)}X_n=R_n^9.
\]
\end{definition}

``Positive integral'' means a finite direct sum of irreducible characters
with coefficients in $\mathbb Z_{\geq0}$.  It does not mean that the lift is
unique, nor that its consecutive difference $X_n-X_{n-1}$ is positive.  It
also does not include Bars's separate higher-multiplicity scheme or arbitrary
additional complete nonperturbative multiplets.  The protected project data
through level 10 provide examples of such lifts; only levels 1--5 were
displayed by Bars himself.

\section{The HFT covariant spectrum}

\subsection{Characters and plethystic exponentials}

Hanany, Forcella and Troost (HFT) supplied the all-order perturbative input
which Bars did not have: a generating function covariant under the massive
little group $\Spin(9)$ \cite{HFT}.  We write
\[
   [a_1,a_2,a_3,a_4]_9
\]
for the irreducible $B_4=\Spin(9)$ character with those Dynkin labels and set
\[
   V_9=[1,0,0,0]_9,
   \qquad S_9=[0,0,0,1]_9.
\]
On the common $\Spin(8)$ torus the physical transverse vector is represented
by the virtual $B_4$ character
\[
   T=V_9-\one,
\]
because $V_9$ restricts as $8_v\oplus\one$.

For a character-valued series $f$, the bosonic and fermionic plethystic
exponentials are
\[
   \PE[f]=\exp\!\left(\sum_{k\geq1}\frac{f(q^k,y^k)}{k}\right),
\]
\[
   \PEF[f]=\exp\!\left(\sum_{k\geq1}
        \frac{(-1)^{k+1}f(q^k,y^k)}{k}\right).
\]
They are compact devices for taking all symmetric, respectively exterior,
powers of oscillator representations.  The variables collectively denoted
by $y$ are weight fugacities.

\subsection{The HFT formula}

Define the transverse bosonic product
\[
   Z_B(q)=\PE\!\left[\frac{q}{1-q}T\right]
\]
and the fermionic product with a fermion-mode fugacity $f$,
\[
   \mathcal Z_F(q;f)
   =\PEF\!\left[\frac{f}{1-q}T\right].
\]
The GSO-projected Neveu--Schwarz series is
\[
   Z_{NS}(q)=\frac{1}{2\sqrt q}
   \left(\mathcal Z_F(q;\sqrt q)
        -\mathcal Z_F(q;-\sqrt q)\right),
\]
while the Ramond oscillator factor entering the covariant expression is
\[
   Z_R(q)=\mathcal Z_F(q;q).
\]
With a purely massless chirality term denoted by $\Delta_0$, HFT's covariant
left-moving partition function can be written
\begin{equation}
\label{eq:hft-left}
   Z_{\mathrm{Left}}(q)
   =\Delta_0+Z_B(q)
      \left(Z_{NS}(q)+\frac12 S_9 Z_R(q)\right).
\end{equation}
The apparently fractional factor does not produce fractional multiplicities:
the plethystic expansion supplies the compensating parity factor.

Let $Z_0$ be the complete massless contribution.  The massive-primary series
is HFT's quotient
\begin{equation}
\label{eq:hft-massive}
   Z_m(q)=\frac{Z_{\mathrm{Left}}(q)-Z_0}{qZ_Q}.
\end{equation}
In project notation,
\begin{equation}
\label{eq:R-series}
   R(q)=Z_m(q)=\sum_{d\geq0}r_{d+1}^9q^d.
\end{equation}
Thus coefficient degree $d=0$ corresponds to physical massive level one.
Equation \eqref{eq:hft-massive} is an exact all-orders character formula, but
positivity of its irreducible $\Spin(9)$ coefficients is not manifest from
the quotient.  HFT themselves identified proving positivity as a challenge.
For the Bars problem, $R(q)$ supplies the single-level characters which are
then cumulatively summed as in \eqref{eq:bars-cumulative}.

\subsection{Removing the leading symmetric-vector tower}

HFT also observed the exact factorization
\begin{equation}
\label{eq:hft-regge-factor}
   R(q)=\PE(qV_9)F(q),
   \qquad
   F(q)=\sum_{d\geq0}F_dq^d.
\end{equation}
The coefficient of the first factor is
\[
   [q^a]\PE(qV_9)
   =\operatorname{Sym}^{a}(V_9)
   =\bigoplus_{j=0}^{\lfloor a/2\rfloor}[a-2j,0,0,0]_9.
\]
This factor carries the leading Regge tower.  The coefficients $F_d$ are
\emph{virtual} characters: irreducible coefficients may be negative because
we have algebraically removed nine level-one modes even though only eight are
physical transverse oscillators.  Therefore an argument about $F$ must
respect cancellations; $F_d$ is not the character of a Hilbert space.

Equation \eqref{eq:hft-regge-factor} is the entry point of the final proof,
but several layers of the lifting problem had to be understood before the
right support statement for $F$ became visible.

\section{Branching packets and positive lifts}

\subsection{The multiplicity-free branching rule}

Write a dominant $D_5=\Spin(10)$ highest weight in orthogonal coordinates as
\[
   \lambda_1\geq\lambda_2\geq\lambda_3\geq\lambda_4
   \geq|\lambda_5|
\]
and a dominant $B_4=\Spin(9)$ weight as
\[
   \mu_1\geq\mu_2\geq\mu_3\geq\mu_4\geq0.
\]
The restriction $D_5\downarrow B_4$ is multiplicity-free.  The child
$V_\mu$ occurs in $V_\lambda$ precisely when the coordinates interlace:
\begin{equation}
\label{eq:interlacing}
 \lambda_1\geq\mu_1\geq\lambda_2\geq\mu_2
 \geq\lambda_3\geq\mu_3\geq\lambda_4
 \geq\mu_4\geq|\lambda_5|.
\end{equation}

\begin{definition}[Branching packet]
The packet of a $D_5$ irrep $V_\lambda$ is the $B_4$ character
\[
   P_\lambda=\Res^{D_5}_{B_4}V_\lambda,
\]
namely the sum of all interlacing children in
\eqref{eq:interlacing}.  Its multiplicity vector has only zeroes and ones.
\end{definition}

A positive integral lift of a target multiplicity vector $m$ is exactly a
nonnegative integral packet decomposition
\[
   m=Ax,
   \qquad x\in\mathbb Z_{\geq0}^{p},
\]
where the columns of $A$ are the permitted packets.  This converts Bars's
question into a structured integer decomposition problem.

For a top $B_4$ Dynkin label $\mu=[\ell,m,r,s]_9$, the unoriented family of
parents with that same top child is
\begin{equation}
\label{eq:parent-family}
   [\ell,m,r,k,s-k]_{10},
   \qquad 0\leq k\leq\left\lfloor\frac{s}{2}\right\rfloor.
\end{equation}
Outer-conjugate $D_5$ parents restrict identically, so they may be collapsed
when the question is existence rather than chirality selection.

\subsection{A triangular virtual section}

Choose the $k=0$ parent
\[
   p_0(\ell,m,r,s)=[\ell,m,r,0,s]_{10}
\]
for every $B_4$ top weight.  Ordered by dominance, the resulting packet
matrix $U$ is triangular with diagonal entries one.  Hence
\[
   \det U=1
\]
in every finite downward-closed domain, and any integral target has a unique
integral \emph{virtual} lift on this section:
\[
   c=U^{-1}m.
\]
Some coefficients $c_\mu$ can be negative.  The role of alternative parents
is precisely to repair those negative coefficients while preserving the
restriction.

\begin{definition}[Restriction kernel and positive fibre]
The restriction kernel consists of virtual $D_5$ characters $K$ satisfying
$\Res K=0$.  If $s(m)$ is any virtual section lift, every other lift is
\[
   s(m)+K.
\]
The \emph{positive fibre} over $m$ is the set of those kernel corrections for
which every final $D_5$ coefficient is a nonnegative integer.
\end{definition}

This language separates existence from selection.  Bars III.B asks whether
the positive fibre is nonempty.  It does not ask for a preferred point in
that fibre.

\section{The trajectory--kernel programme}

\subsection{Step 1a: Regge coordinates and shifts}

At each physical HFT level $n\geq1$, write
\[
 a_n(\ell,m,r,s)
 =\operatorname{mult}_{r_n^9}[\ell,m,r,s]_9.
\]
Thus a five-tuple $(n,\ell,m,r,s)$ means ``level $n$, with $B_4$ highest
weight $[\ell,m,r,s]_9$''.  The first Dynkin coordinate $\ell$ measures the
leading symmetric-tensor direction.  The basic Regge shift is
\begin{equation}
\label{eq:regge-shift}
   (n,\ell,m,r,s)\longmapsto(n+1,\ell+1,m,r,s).
\end{equation}
A \emph{Regge column} fixes $(m,r,s)$.  The elementary orbit invariant is
therefore
\[
   \rho=n-\ell.
\]
At fixed $(m,r,s)$, a Regge orbit is therefore specified by $\rho$.  Nothing
about the shift itself singles out the further subtraction of $3m$; that
choice becomes natural only after the support geometry is introduced.

\subsection{Step 1b: HFT support and the GMSZ chambers}

Through the tested finite window, the HFT data are accurately organized by
the empirical support rule of Georgoudis, Minahan, Str\"om and Zoumis (GMSZ)
\cite{GMSZ}.  Their rule is stated for $n>6$.  Define the support cost
\[
 c(r,s)=\max\{6r+5s-6,\ 6r+4s-2,\ 6r+3s-1,\ 2s+1\}.
\]
At fixed $(n,m,r,s)$, the largest allowed leading label is
\[
 \ell_{\max}=n-3m-c(r,s),
\]
apart from the stated isolated gap $[n-2,0,0,0]_9$.  Thus the outer support
condition is
\[
 0\leq\ell\leq\ell_{\max}.
\]

This is what motivates the support-adapted reduced depth
\[
   \delta=n-\ell-3m=\rho-3m.
\]
The subtraction of $3m$ is not required by Regge invariance.  It absorbs the
three units of level carried by one unit of the second Dynkin label in the
GMSZ support rule.  The support condition becomes simply
\[
 \delta\geq c(r,s),
\]
and the literal distance below the leading edge is
\[
 \ell_{\max}-\ell=\delta-c(r,s).
\]
Thus $\rho=n-\ell$ is the basic Regge invariant, whereas $\delta$ is the
same invariant recentered so that the $m$-dependent support cost has already
been removed.

Some chamber terminology will be used repeatedly.  A supported target with
$\ell\geq1$ has the unit Regge predecessor
\[
 (n-1,\ell-1,m,r,s).
\]
In the GMSZ range, its \emph{upper $\ell$-chamber} is the high-$\ell$ part of
the support for which that predecessor is present and
\[
 \ell>\left\lfloor\frac{\ell_{\max}}2\right\rfloor.
\]
Here ``upper'' means closer to the largest allowed first Dynkin label; it
does not mean higher string level.  The multiplicity in this chamber is
transported unchanged from the predecessor.  The lattice wall at the edge
of this region is the \emph{transport boundary}.  The \emph{outer support
boundary} is instead the edge of the allowed spectrum itself, such as
$\ell=\ell_{\max}$ at fixed $(n,m,r,s)$.  A third object, the branching
boundary of an interlacing packet, will appear in Step~2.  These boundaries
are related manifestations of failed naive translation, but they are not the
same mathematical locus.

We can now state the three transport rules precisely in our conventions.
Fix a supported target $[\ell,m,r,s]_9$ at level $n>6$, and let
$\ell_{\max}$ always mean the target value defined above.  A proposed
equality is admissible only when its lower-level source exists in the HFT
support and the indicated positive integer $k$ is within range.

The \emph{$\ell$-transport} is
\[
 a_n(\ell,m,r,s)=a_{n-k}(\ell-k,m,r,s),
 \qquad 1\leq k\leq\ell,
\]
with chamber inequality
\[
 \ell>\left\lfloor\frac{\ell_{\max}}2\right\rfloor
       +\left\lfloor\frac{k}{2}\right\rfloor.
\]
For $k=1$ this reduces to the upper-$\ell$-chamber condition already given.

The \emph{$m$-transport} is
\[
 a_n(\ell,m,r,s)=a_{n-3k}(\ell,m-k,r,s),
 \qquad 1\leq k\leq m,
\]
with chamber inequality
\[
 \ell>\ell_{\max}-m+k-1.
\]
The factor three in the level shift is the same factor that motivated the
coordinate $\delta=n-\ell-3m$.

The \emph{$r$-transport} is
\[
 a_n(\ell,m,r,s)=a_{n-6k}(\ell,m,r-k,s),
 \qquad 1\leq k\leq r,
\]
with chamber inequality
\[
 \ell>\ell_{\max}-r+k-1+\epsilon(s),
\]
where the spinor-label correction is
\[
 \epsilon(s)=
 \begin{cases}
  2, & s=0,\\
  1, & 1\leq s\leq3,\\
  0, & s\geq4.
 \end{cases}
\]
All three inequalities are strict inequalities between integers.  In the
deterministic TK5 chamber decomposition, the families are tried in the order
$\ell$, $m$, $r$, using the largest admissible $k$ in the first applicable
family.  The finite audit also tested every admissible $k$, rather than only
that selected representative.

The \emph{residual core} is the set of supported labels for which none of
these admissible transports applies.  It is therefore the part of the
spectrum whose multiplicity is not supplied by the transport identities.  In
particular, it is not merely a synonym for the complement of the upper
$\ell$-chamber: some labels below that chamber can still be transported in
the $m$ or $r$ direction.

\subsection{Step 1c: what TK1 and TK5 established}

TK1 asked, before imposing any chamber condition, whether HFT multiplicities
were simply constant along all Regge orbits.

They are not globally constant.  Through level 10, 80 of 134 adjacent tests
matched and 54 failed.  TK5 later showed that the tested upper $\ell$-chamber
contains no such failures.  In the later level-$7$--$11$ audit, all 95
admissible adjacent comparisons inside that chamber matched.  The failures
lie outside it, where crossing the transport boundary and encountering
genuinely nontransported core data cannot be ignored.

This chamber decomposition explained the TK1 mismatches, but its status was
still empirical and finite-level.  It prevented a false all-orders
conjecture based only on leading trajectories.  The later Pieri argument
derives directly from the exact HFT factorization the portion of stable
transport actually needed for the all-orders support theorem.

\subsection{Step 2: boundary-corrected kernel shifts}

Let $\sigma_D$ and $\sigma_B$ increase the first $D_5$ and $B_4$ Dynkin
labels respectively.  Interlacing gives the exact identity
\begin{equation}
\label{eq:branch-shift-defect}
   \Res(\sigma_DV_\lambda)
   =\sigma_B\Res(V_\lambda)+\partial V_\lambda,
\end{equation}
where $\partial V_\lambda$ is the nonnegative, multiplicity-free slice
\[
   \mu_1=\lambda_2.
\]
This slice is the \emph{branching boundary}: it is the face of the
$D_5\downarrow B_4$ interlacing packet left behind by the shift.  It is not
the transport boundary just defined, although both phenomena express a
failure of a naive shift to act everywhere without correction.  The slice is
independent of the shifted first label.  Thus raw shifting does not commute
with restriction, but its failure is an explicitly known lower face of the
packet.

If $s_{\mathrm{ch}}$ denotes the triangular all-chiral section, define the
corrected shift on a kernel element by
\begin{equation}
\label{eq:corrected-kernel-shift}
   T(K)=\sigma_DK-s_{\mathrm{ch}}(\Res(\sigma_DK)).
\end{equation}
Then $\Res T(K)=0$.  TK2 proved more: for every natural finite generator,
$T(K)$ is exactly the next generator along its Regge trajectory.  This is
why the programme is called \emph{trajectory--kernel}: the leading shift acts
naturally only after the restriction-kernel boundary correction is included.

The first important seed is the six-term relation
\[
 [0,0,0,1,1]_{10}-[0,0,0,0,2]_{10}
 -[0,0,1,0,0]_{10}+[0,1,0,0,0]_{10}
 -[1,0,0,0,0]_{10}+[0,0,0,0,0]_{10},
\]
whose restriction vanishes.  In dimension notation it is
\[
   210-126-120+45-10+1.
\]
This is the local algebra behind the first $126/210$ parent ambiguity.

\subsection{Step 3: positive fibres and the failure of uniqueness}

TK3 used the natural kernel basis to parameterize the complete positive
fibres in the level-10 HFT closure.  The number of positive lifts at levels
1 through 10 was
\[
   1,1,1,1,1,1,2,8,72,6057.
\]
The first compulsory correction occurs at level 7 and uses the six-term seed
above.  The result establishes two logically different facts:
\begin{itemize}
  \item the tested cumulative HFT targets do have positive integral lifts;
  \item positivity alone very quickly ceases to choose a unique lift.
\end{itemize}

TK4 and TK5 tested possible selection principles: minimum kernel population,
minimum support, minimum corrected-shift defect and sources confined to an
empirical residual chamber.  These organize the finite solutions but do not
produce a unique cutoff-stable rule.  In particular, a lift optimized at one
horizon may have to be abandoned to reach the optimum at the next horizon.
This redirected the project from selecting a favourite lift to proving that
the positive fibre is nonempty by intrinsic inequalities.

\section{From a positive fibre to facet inequalities}

\subsection{The global ring kernel}

Let $16_+$ and $16_-$ be the two half-spinor characters of $D_5$ and put
\[
   J=16_+-16_-.
\]
The representation-ring restriction is surjective and has the exact
principal kernel
\begin{equation}
\label{eq:principal-kernel}
   \ker\!\left(R(\Spin(10))\longrightarrow R(\Spin(9))\right)=(J).
\end{equation}
Thus every virtual ambiguity is a multiple of one global chirality
difference.  This elegant ring statement does not solve positivity:
multiplication by $J$ introduces both signs.  The positive problem lives in
the representation \emph{semiring}, where subtraction is unavailable.

\subsection{Cone, semigroup, facets and Hilbert basis}

For a finite packet matrix $A$, define
\[
   C=A\mathbb R_{\geq0}^{p},
   \qquad
   S=A\mathbb Z_{\geq0}^{p}.
\]
The cone $C$ records rational positive packet combinations; the semigroup
$S$ records genuine integral combinations.  An integral target can lie in
$C$ but outside $S$, producing an \emph{integrality hole}.

A \emph{facet} is a codimension-one face of $C$.  Equivalently it is given by
a primitive linear inequality
\[
   y^Tm\geq0
\]
valid on every packet column.  Membership in a full-dimensional rational
cone is equivalent to satisfying all its facet inequalities.  A
\emph{Hilbert basis} is the minimal finite set of lattice points which
generates all lattice points in the cone by nonnegative integer sums.  If the
packet columns themselves form that Hilbert basis, then
\[
   S=C\cap\mathbb Z^d,
\]
and the semigroup is called \emph{saturated} or \emph{normal}.

TK6B.1 found, in the complete level-11 closure,
\[
   A\in\mathbb Z^{73\times84},
   \qquad \operatorname{rank}A=73,
\]
with 2508 facets.  The matrix is not totally unimodular, but its 84 packet
columns are nevertheless exactly its Hilbert basis.  Hence, in that finite
domain, the facet inequalities are both necessary and sufficient for an
integral positive lift.  This was decisive conceptually, although it was not
yet an all-orders saturation theorem.

\subsection{Gale transforms and positive circuits}

Directly listing thousands of facet inequalities obscures their origin.  Let
\[
   K\in\mathbb Z^{p\times k}
\]
have as columns an integral basis of $\ker A$, so $AK=0$.  The row
vectors $g_j\in\mathbb Z^k$ of $K$ are a \emph{Gale transform} of the packet
configuration.  Informally, the packet columns describe what is generated;
the Gale vectors describe all linear dependencies among those generators.

If $y$ is a facet functional, its slack on the packet columns is
\[
   x=A^Ty.
\]
It satisfies
\[
   x\geq0,
   \qquad K^Tx=0.
\]
Conversely, because $A$ has full row rank, every extreme ray of this
nonnegative relation cone comes from a facet.  A \emph{positive Gale circuit}
is a minimally supported relation
\[
   \sum_j x_jg_j=0,
   \qquad x_j\geq0,
\]
with primitive integer coefficients.  The word ``circuit'' is borrowed from
matroid theory: it means that deleting any supported row destroys the
dependence.

In triangular coordinates $c=U^{-1}m$, the corresponding liftability
condition becomes the weighted Hall-type inequality
\begin{equation}
\label{eq:weighted-circuit}
   \sum_{j\in\operatorname{Ref}}x_jc_j\geq0.
\end{equation}
At level 11, 2508 facets compressed to only 40 circuit types.  At level 13,
new shell collisions produced genuine coefficient-two circuits.  Therefore
the correct all-orders object is a growing family of \emph{weighted}
circuits, not a fixed list of zero-one inequalities.

\section{Shell identities and the bridge back to HFT}

\subsection{Universal local kernel relations}

The finite kernel families eventually proliferate, but their local origin is
uniform.  For the parent family \eqref{eq:parent-family}, let
\[
   P_k=[\ell,m,r,k,s-k]_{10}.
\]
The difference of two consecutive parent packets is a single interlacing
boundary slice.  Inclusion--exclusion on that slice gives four all-chiral
packets $Q_{1,k},\ldots,Q_{4,k}$ and the universal six-term identity
\begin{equation}
\label{eq:shell-identity}
   \Res\bigl(P_k-P_{k-1}-Q_{1,k}+Q_{2,k}-Q_{3,k}+Q_{4,k}\bigr)=0.
\end{equation}
For example, the first split is
\begin{multline}
 K_{\ell,m,r,s}^{(1)}=
 [\ell,m,r,1,s-1]-[\ell,m,r,0,s]
 -[\ell,m,r+1,0,s-2]\\
 +[\ell,m+r+1,0,0,s-2]
 +[m,r,0,0,s-2]-[\ell+m+1,r,0,0,s-2].
\end{multline}
All labels in this display are $D_5$ labels.  The identity follows directly
from the interlacing rule and is valid for arbitrary nonnegative labels.
Summing consecutive shells telescopes the $k$th parent back to the all-chiral
parent with exactly $4k+2$ unit-coefficient terms.

The important point is where level dependence has gone.  The shell formula
is universal; a string-level cutoff merely chooses a finite induced
collision hypergraph of these local relations.  Its primitive positive
circuits are the facet inequalities that the HFT target must satisfy.

\subsection{The local M\"obius inverse}

The triangular section itself also has a cutoff-independent formula.  Let
\[
   M(\ell,m,r,s)
\]
be the multiplicity of $[\ell,m,r,s]_9$ in any finite $B_4$ character,
extended by zero outside the dominant chamber.  The coefficient of the
all-chiral parent $[\ell,m,r,0,s]_{10}$ is
\begin{equation}
\label{eq:mobius}
\begin{split}
 c_{\ell,m,r,s}
 ={}&\sum_{\epsilon_1,\epsilon_2,\epsilon_3\in\{0,1\}}
 (-1)^{\epsilon_1+\epsilon_2+\epsilon_3}\\
 &M(\ell+\epsilon_1-\epsilon_2,
    m+\epsilon_2-\epsilon_3,
    r+\epsilon_3,s).
\end{split}
\end{equation}
This is a three-fold finite difference, or M\"obius inversion, of the
interlacing sums.  ``M\"obius'' here means inclusion--exclusion on a finite
partially ordered set; no number-theoretic M\"obius function is involved.

Substituting \eqref{eq:mobius} into
\eqref{eq:weighted-circuit} turns every branching-cone facet into a finite
signed inequality among HFT multiplicities.  This is the essential bridge:
an abstract $\Spin(10)$ lifting obstruction becomes an explicit local
functional of the known $\Spin(9)$ generating series.

\subsection{Why higher-level enumeration cannot finish the proof}

The full infinite shell system has no nonzero finitely supported positive
circuit.  The reason is that one boundary term in
\eqref{eq:shell-identity} forgets the leading label $\ell$, so a putative
finite circuit can always be tested against a sufficiently distant shell and
forced negative.  Finite-horizon facets are therefore moving
\emph{boundary-tail} objects; their support changes as the horizon grows.

Some tails telescope under \eqref{eq:mobius} to a bounded number of endpoint
multiplicities, but mixed circuits can have growing pullback support.  This is
the precise answer to the question ``what prevents an all-orders proof from
being just a check of higher and higher levels?''  One needs a structural
transport theorem for the moving tails.  That requirement led back to the
HFT symmetric-vector factorization.

\section{Stable Pieri translation and the remaining support theorem}

\subsection{What ``stable Pieri translation'' means}

Let $V_\mu$ and $V_\lambda$ be irreducible $B_n$ representations, expressed
in orthogonal coordinates, and let
\[
   H_k=V_{k\omega_1}
\]
be the traceless symmetric, or harmonic one-row, representation.  A Pieri
rule describes the irreducible summands obtained by tensoring an arbitrary
representation with such a one-row representation.

\begin{theorem}[Stable Pieri translation]
If the prospective target is at least as far out in its first orthogonal
coordinate as the source,
\[
   \lambda_1\geq\mu_1,
\]
then
\begin{equation}
\label{eq:stable-pieri-harmonic}
 [V_\lambda:V_\mu\otimes H_k]
 =
 [V_{\lambda+\omega_1}:V_\mu\otimes H_{k+1}]
\end{equation}
for every $k\geq0$.  Consequently
\begin{equation}
\label{eq:stable-pieri-symmetric}
 [V_\lambda:V_\mu\otimes\operatorname{Sym}^{a}(V)]
 =
 [V_{\lambda+\omega_1}:
  V_\mu\otimes\operatorname{Sym}^{a+1}(V)].
\end{equation}
\end{theorem}

The brackets denote tensor-product multiplicity.  In $B_4$,
$\lambda\mapsto\lambda+\omega_1$ is exactly
$[\ell,m,r,s]\mapsto[\ell+1,m,r,s]$.  Thus the theorem says that, once the
first row is longer than anything carried by the fixed source $V_\mu$, adding
one symmetric-vector degree simply adds one box to that first row.  The
decomposition has entered a stable regime.

The proof uses Racah--Speiser.  Every Weyl-group term which moves the first
coordinate would require a weight of $V_\mu$ whose absolute coordinate is
larger than $\mu_1$, so it vanishes.  The stabilizer of the first coordinate
is unchanged when one adds $\omega_1$ to the target and one unit to the
harmonic degree.  This works equally for integral tensor weights and
half-integral spinor weights; stable Pieri translation itself is not a
chirality obstruction.

\subsection{Applying it to the HFT factorization}

From \eqref{eq:hft-regge-factor}, the coefficient of $q^N$ is
\[
   [q^N]R(q)
   =\sum_{d=0}^{N}\operatorname{Sym}^{N-d}(V_9)\otimes F_d.
\]
For a fixed $V_\mu$ in $F_d$, equation
\eqref{eq:stable-pieri-symmetric} transports its contribution from one
physical level to the next provided the target satisfies
$\lambda_1\geq\mu_1$.  The sign of the coefficient of $V_\mu$ in the virtual
character $F_d$ is irrelevant: the equality holds contribution by
contribution before the signed terms are added.

Two possible failures remain:
\begin{enumerate}
  \item a source $V_\mu$ might couple to an \emph{unstable} target with
        $\lambda_1<\mu_1$;
  \item the new top-degree term $F_N$, which has no predecessor, might itself
        hit the translated target.
\end{enumerate}
Both are excluded if every remainder label appears late enough.

\subsection{The Pieri hazard}

Define the exterior-chamber cost
\[
   E(\ell,m,r,s)=2\ell+3m+c(r,s),
\]
where $c(r,s)$ is the support cost introduced above.  For a source $\mu$,
define its harmonic hazard
\[
 H(\mu)=
 \max_{\substack{k\geq0,\ V_\lambda\subset V_\mu\otimes H_k\\
                  \lambda_1<\mu_1}}
 \bigl(E(\lambda)-k\bigr).
\]
This maximum is finite.  An unstable target can occur only for
\[
   k<2\mu_1,
\]
so no unbounded tensor-product search is hidden in the definition.  Put
\begin{equation}
\label{eq:pieri-threshold}
   K(\mu)=\max\bigl(E(\mu)-1,H(\mu)\bigr).
\end{equation}
Then the single support implication
\begin{equation}
\label{eq:needed-support}
   [V_\mu]F_d\neq0
   \quad\Longrightarrow\quad
   d\geq K(\mu)
\end{equation}
simultaneously excludes all unstable couplings and the unpaired top-degree
term.

The hazard is not an opaque maximization.  Orthogonal Pieri interlacing
\cite{Okada} gives a closed piecewise-affine formula: on $\ell=0$ it is
simply $E-1$, while for $\ell\geq1$ three additional affine walls appear.
The next section packages those pieces into a convenient seven-wall function
$K_7$.

\subsection{What the present proof accomplishes}

Before TK6B.5f, implication \eqref{eq:needed-support} had been checked only
through a finite HFT window.  Checking more degrees could never establish it.
The remaining task was to derive it from the infinite oscillator product in
\eqref{eq:hft-left}.  The rest of this report gives that derivation:
\begin{enumerate}
  \item encode $K$ by seven affine walls;
  \item rewrite the HFT quotient on the common $D_4$ torus;
  \item telescope the Green--Schwarz product using a finite triality identity;
  \item prove that the infinite high-mode tail is automatically safe;
  \item check the fixed low-mode residue exactly.
\end{enumerate}
This closes one major analytic premise of the trajectory--kernel programme.
It does not yet classify all moving circuits or prove all-level integral
packetization.

\subsection{Programme at a glance}

The route can be summarized without the finite computational detail as
follows.  A ``complete'' entry in the final column means complete for the
mathematical statement indicated there, not completion of the Bars problem.

\begin{longtable}{p{0.12\textwidth}p{0.35\textwidth}p{0.43\textwidth}}
\caption{Conceptual progression of the trajectory--kernel programme.}
\label{tab:programme-roadmap}\\
\toprule
stage & new object or theorem & consequence \\
\midrule
\endfirsthead
\toprule
stage & new object or theorem & consequence \\
\midrule
\endhead
Bars & cumulative target $R_n^9$ & formulate the minimal positive
$D_5$-lift existence problem \\
HFT & covariant series $R(q)$ & replace order-by-order oscillator
construction by an exact all-orders $B_4$ generating function \\
TK1 & Regge columns and reduced depth $\delta$ & raw multiplicity shifts are
useful but not globally valid \\
TK2 & boundary-corrected kernel shift & exact transport acts on restriction
relations rather than on raw packets \\
TK3--5 & positive fibres and chamber/core sources & establish finite
existence and strong nonuniqueness; simple minimization is not canonical \\
TK6a & principal kernel ideal $(16_+-16_-)$ & settle the global virtual
kernel while isolating positivity as a semiring problem \\
TK6B.1 & packet cone, semigroup and Hilbert basis & reduce finite integral
liftability to facet inequalities in the level-11 closure \\
TK6B.2--4 & Gale circuits and universal shells & replace thousands of facets
by primitive weighted relations in a local collision hypergraph \\
TK6B.5b/c & M\"obius pullback and boundary tails & turn facets into local HFT
multiplicity inequalities and explain why circuit lists move with the cutoff
\\
TK6B.5d/e & stable Pieri translation and hazard $K$ & reduce the relevant
transport inequalities to the single support implication
$[V_\mu]F_d\neq0\Rightarrow d\geq K(\mu)$ \\
TK6B.5f & Green--Schwarz telescope and finite low-mode core & prove that
support implication at every degree \\
\bottomrule
\end{longtable}

\section{The seven-wall degree}

\subsection{$B_4$ labels and orthogonal coordinates}

We write a dominant $B_4$ Dynkin label as
\[
   \mu=(\ell,m,r,s).
\]
Its orthogonal coordinates are
\begin{align*}
 x_1&=\ell+m+r+\frac{s}{2}, &
 x_2&=m+r+\frac{s}{2},\\
 x_3&=r+\frac{s}{2}, &
 x_4&=\frac{s}{2},
\end{align*}
with $x_1\geq x_2\geq x_3\geq x_4\geq0$.  Tensorial representations have
integral $x_i$, while spinorial representations have half-integral $x_i$.

The original Pieri threshold $K(\mu)$ is piecewise affine.  It is convenient
to use a slightly stronger uniform function $\Kseven$, which agrees with $K$
away from the wall $\ell=0$ and may be stronger on that wall.

\begin{definition}[Uniform seven-wall degree]
For $\mu=(\ell,m,r,s)$, define $\Kseven(\mu)$ as the maximum of
\begin{align*}
 W_1&=2\ell+3m+6r+5s-7,\\
 W_2&=2\ell+3m+6r+4s-3,\\
 W_3&=2\ell+3m+6r+3s-2,\\
 W_4&=2\ell+3m+2s,\\
 W_5&=2\ell+5m+6r+3s-4,\\
 W_6&=2\ell+5m+7r+4s-5,\\
 W_7&=2\ell+5m+9r+5s-9.
\end{align*}
\end{definition}

The first four walls are the exterior-chamber degree $E(\mu)-1$; the last
three are the extra unstable-Pieri hazards.  Since
\[
   \Kseven(\mu)\geq K(\mu),
\]
it is enough to prove the stronger implication
\[
   [V_\mu]F_d\neq0\quad\Longrightarrow\quad d\geq\Kseven(\mu).
\]

For later use, Table~\ref{tab:orthogonal-walls} gives the same seven walls in
orthogonal coordinates.

\begin{table}[ht]
\centering
\caption{The seven walls in orthogonal coordinates.}
\label{tab:orthogonal-walls}
\begin{tabular}{ccl}
\toprule
wall & coefficient vector & affine expression \\
\midrule
$W_1$ & $(2,1,3,4)$ & $2x_1+x_2+3x_3+4x_4-7$ \\
$W_2$ & $(2,1,3,2)$ & $2x_1+x_2+3x_3+2x_4-3$ \\
$W_3$ & $(2,1,3,0)$ & $2x_1+x_2+3x_3-2$ \\
$W_4$ & $(2,1,-3,4)$ & $2x_1+x_2-3x_3+4x_4$ \\
$W_5$ & $(2,3,1,0)$ & $2x_1+3x_2+x_3-4$ \\
$W_6$ & $(2,3,2,1)$ & $2x_1+3x_2+2x_3+x_4-5$ \\
$W_7$ & $(2,3,4,1)$ & $2x_1+3x_2+4x_3+x_4-9$ \\
\bottomrule
\end{tabular}
\end{table}

\section{Changing to the Green--Schwarz oscillator product}

\subsection{Why $D_4$ is the useful intermediate language}

The covariant HFT series is naturally organized as a $B_4=\Spin(9)$
character.  The oscillator product is simpler after restriction to the common
maximal torus of $D_4=\Spin(8)$.  Let
\[
   V=8_v,
   \qquad S=8_s,
   \qquad C=8_c,
\]
where $S$ is the chirality chosen for the Green--Schwarz fermionic
oscillators and $C$ is the other half-spinor.  Exchanging $S$ and $C$ is an
outer automorphism of $D_4$; no physical chirality choice for a
$\Spin(10)$ lift is being made at this stage.

Define the finite exterior polynomials
\[
   D(u)=\sum_{j=0}^{8}(-u)^j\Lambda^jV,
   \qquad
   E(u)=\sum_{j=0}^{8}u^j\Lambda^jS,
\]
and set
\[
   H=V+S.
\]
The inverse $D(u)^{-1}$ is the symmetric-power series of $V$.

\subsection{The finite triality identity}

\begin{lemma}[One-mode divisibility]
In the exact representation ring of $D_4$,
\[
   E(u)-D(u)=uH B(u),
\]
where
\[
   B(u)=1+(C-\one)u^2+(V-S)u^3+(C-\one)u^4+u^6.
\]
\end{lemma}

\begin{proof}
Both sides are polynomials of degree at most eight.  Expand the exterior
powers of $8_v$ and $8_s$.  Triality identifies the relevant vector and
spinor exterior powers, while Hodge duality pairs degrees $j$ and $8-j$.
After subtracting, the constant terms cancel and every remaining coefficient
is divisible by $H=V+S$.  The quotients, in increasing degree, are
\[
   \one,\quad 0,\quad C-\one,\quad V-S,
   \quad C-\one,\quad 0,\quad\one.
\]
This is exactly $B(u)$.  The implementation verifies all nine coefficients
as identities of Weyl characters.
\end{proof}

The identity is a finite representation-ring form of the generalized Jacobi
identity relating RNS and Green--Schwarz oscillator descriptions.  Its main
advantage is that it performs the ground-multiplet division without expanding
an infinite virtual Koszul sum.

\subsection{Telescoping the infinite product}

Put
\[
   A(u)=\frac{E(u)}{D(u)}.
\]
The lemma gives
\[
   A(u)-1=\frac{uH B(u)}{D(u)}.
\]
Let
\[
   P(q)=\prod_{m\geq1}A(q^m)
\]
be the Green--Schwarz oscillator product.  For a finite product one has the
elementary telescoping identity
\[
   \prod_{m=1}^{N}A_m-1
   =\sum_{n=1}^{N}(A_n-1)\prod_{m=1}^{n-1}A_m.
\]
Substituting $A_n=A(q^n)$ and using the one-mode identity gives, coefficient
by coefficient in $q$,
\[
 \frac{P(q)-1}{qH}
 =\sum_{n\geq1}q^{n-1}\frac{B(q^n)}{D(q^n)}
   \prod_{m=1}^{n-1}A(q^m).
\]
There is no convergence issue: the coefficient of $q^d$ sees only finitely
many modes.

The $B_4$ vector restricts to $\one+V$ on $D_4$.  Removing its level-one
symmetric tower therefore multiplies by
\[
   \PE(-q(\one+V))=(1-q)D(q).
\]
We obtain the central exact identity.

\begin{proposition}[Green--Schwarz telescope for the remainder]
On the common $D_4/B_4$ torus,
\begin{equation}
\label{eq:telescope}
 F(q)=(1-q)\sum_{n\geq1}q^{n-1}B(q^n)
 \frac{\displaystyle\prod_{m=1}^{n-1}E(q^m)}
      {\displaystyle\prod_{m=2}^{n}D(q^m)}.
\end{equation}
\end{proposition}

This is the decisive change of viewpoint.  The universal multiplet has been
divided out exactly, with the required cancellations already incorporated.
Every summand has a distinguished first nontrivial mode $n$, which we call
its \emph{pivot}.

\section{Returning from $D_4$ to $B_4$}

\subsection{The finite Weyl-numerator map}

Let $A_\rho$ denote a Weyl alternant.  Since the roots of $B_4$ not present
in $D_4$ are the short roots $\pm e_i$, the Weyl denominators obey
\[
   \frac{A_{\rho_{B_4}}}{A_{\rho_{D_4}}}
   =\prod_{i=1}^{4}\left(e^{e_i/2}-e^{-e_i/2}\right).
\]
For a $D_4$ irrep with highest weight $\lambda$,
\[
 A_{\rho_{B_4}}\chi^{D_4}_{\lambda}
 =\frac{A_{\rho_{B_4}}}{A_{\rho_{D_4}}}
  A_{\lambda+\rho_{D_4}}.
\]
The right side has at most $192\cdot16$ signed monomials before equal terms
are collected: 192 elements of the $D_4$ Weyl group and 16 sign choices from
the short-root product.  Thus the conversion is finite and independent of
the size of $\lambda$.

\subsection{Chirality folding does not increase the wall degree}

Write a dominant $D_4$ highest weight in orthogonal coordinates as
\[
   \lambda=(a,b,c,d),
   \qquad a\geq b\geq c\geq|d|.
\]
Its chirality fold is the dominant $B_4$ weight
\[
   \widehat\lambda=(a,b,c,|d|).
\]

\begin{lemma}[Componentwise containment]
If a dominant $B_4$ term
$\nu=(y_1,y_2,y_3,y_4)$ occurs in
$A_{\rho_{B_4}}\chi^{D_4}_\lambda$, then
\[
   y_i\leq\widehat\lambda_i
   \qquad (1\leq i\leq4).
\]
\end{lemma}

\begin{proof}
The doubled Weyl-numerator terms have the form
\[
   w(2\lambda+2\rho_{D_4})+\varepsilon,
   \qquad \varepsilon_i\in\{\pm1\}.
\]
Take absolute values and order the four coordinates decreasingly.  There is
only one coordinate originating from each entry of
$\widehat\lambda+\rho_{D_4}$.  Consequently the $i$th order statistic is no
larger than
\[
   \widehat\lambda_i+(\rho_{D_4})_i+\frac12
   =\widehat\lambda_i+(\rho_{B_4})_i.
\]
Subtracting $\rho_{B_4}$ gives the claim.
\end{proof}

Six of the seven walls in Table~\ref{tab:orthogonal-walls} have nonnegative
orthogonal-coordinate coefficients, so componentwise containment immediately
controls them.  Only $W_4$ needs attention.

\begin{lemma}[The exceptional wall]
Every dominant term $\nu$ in the Weyl numerator satisfies
\[
   \Kseven(\nu)\leq\Kseven(\widehat\lambda).
\]
\end{lemma}

\begin{proof}
For the exceptional wall, dominance of $\nu$ gives
\[
\begin{split}
 W_4(\nu)
 &=2y_1+y_2-3y_3+4y_4\\
 &\leq2a+b+|d|,
\end{split}
\]
because $-3y_3+4y_4\leq y_4\leq|d|$.
If $c=|d|$, the last expression is exactly $W_4(\widehat\lambda)$.
If $c>|d|$, the two coordinates lie in the same weight lattice, so
$c\geq|d|+1$.  Hence
\[
   2a+b+|d|\leq2a+b+3c-2=W_3(\widehat\lambda).
\]
The maximum of all seven source walls therefore bounds every target wall.
\end{proof}

This lemma explains precisely what chirality does and does not do here.
Different $D_4$ chiral representatives may contribute differently, but after
folding, none can create a $B_4$ term of larger filtered degree.  No choice of
a $D_5=\Spin(10)$ parent has entered the argument.

\section{Why only finitely many oscillator modes remain}

\subsection{Vector and spinor Lipschitz constants}

The seven walls provide uniform tensor-product estimates.

\begin{lemma}[One-factor degree increase]
Let $\alpha$ be a dominant $D_4$ weight and fold all $D_4$ chiralities before
evaluating $\Kseven$.
\begin{enumerate}
  \item Every constituent of $V_\alpha\otimes V$ has degree at most
        $\Kseven(\widehat\alpha)+4$.
  \item Every constituent of $V_\alpha\otimes S$ has degree at most
        $\Kseven(\widehat\alpha)+5$.
\end{enumerate}
\end{lemma}

\begin{proof}
A vector weight changes one orthogonal coordinate by $\pm1$ before Weyl
dominantization.  The largest absolute coefficient in
Table~\ref{tab:orthogonal-walls} is 4.  A spinor weight changes all four
coordinates by $\pm\tfrac12$.  Half the largest sum of absolute wall
coefficients is
\[
   \frac12(2+1+3+4)=5.
\]
Taking the maximum over walls and Weyl images preserves these bounds.
\end{proof}

Repeated tensoring gives an increase of at most $4r$ for a constituent of
$V^{\otimes r}$ and at most $5j$ for a constituent of $S^{\otimes j}$.
Since $\operatorname{Sym}^rV\subset V^{\otimes r}$ and
$\Lambda^jS\subset S^{\otimes j}$, the same upper bounds apply to the
oscillator factors.

\subsection{The safe tail}

In a bosonic mode $m$, the degree charge of $\operatorname{Sym}^rV$ is $mr$.
Thus every vector mode $m\geq4$ pays at least the maximum possible increase
$4r$.  Similarly, every fermion mode $m\geq5$ pays at least $5j$ for
$\Lambda^jS$.

The polynomial $B(q^n)$ is also harmless for $n\geq3$:
\begin{center}
\begin{tabular}{ccc}
\toprule
power in $B$ & possible non-scalar factor & degree charge for $n\geq3$ \\
\midrule
2 & $C$ & at least 6 \\
3 & $V$ or $S$ & at least 9 \\
4 & $C$ & at least 12 \\
\bottomrule
\end{tabular}
\end{center}
The possible filtered increases are at most 5, 5 and 5 respectively.
Negative coefficients such as $-\one$ or $-S$ do not change a support bound;
they change only signed multiplicities.

We have therefore reduced the entire infinite product to
\begin{itemize}
  \item fermionic modes 1, 2, 3 and 4;
  \item bosonic vector modes 2 and 3;
  \item the exceptional pivot polynomials $B(q)$ and $B(q^2)$.
\end{itemize}
This reduction is uniform in the final string level.

\section{The two low vector modes}

\subsection{Harmonic decomposition and the hazard transform}

Let $H_k$ be the $D_4$ harmonic one-row representation.  Recall
\[
   \operatorname{Sym}^{r}V
   =H_r\oplus H_{r-2}\oplus H_{r-4}\oplus\cdots.
\]
For an upper support estimate, a trace term is easier than its harmonic
label suggests: it pays the oscillator degree of $r$ but has harmonic degree
$k\leq r$.  It is therefore enough to charge $mk$ to $H_k$ in mode $m$.

For a folded source $\alpha$, define the exact mode-two hazard
\[
 H_2(\alpha)=
 \max_{\substack{k\geq0\\V_\beta\subset V_\alpha\otimes H_k}}
 \left(\Kseven(\widehat\beta)-2k\right).
\]
This is finite: after the first-row stable range begins, increasing $k$ only
translates the leading coordinate and does not enlarge the charged degree.

\subsection{Pieri-polytope optimization}

The orthogonal horizontal-strip Pieri rule can be expressed using an
intermediate partition $\xi$.  Both the source-to-$\xi$ skew diagram and the
target-to-$\xi$ skew diagram are horizontal strips.  The number of boxes
removed and added bounds the harmonic degree $k$.  In the spinorial coset,
subtracting the common half-spin shift gives the same interlacing
inequalities.

Each wall is affine in the target rows, while the degree penalty is affine in
the source, intermediate and target rows.  Maximizing over the Pieri polytope
is therefore an elementary endpoint calculation: every variable is put at
one of its two interlacing endpoints according to the sign of its coefficient.
This gives the certified upper closures in
Table~\ref{tab:mode-two}.

\begin{table}[ht]
\centering
\caption{Affine upper closure under vector mode 2.  A row
$(p_1,p_2,p_3,p_4)-c$ denotes
$p_1\ell+p_2m+p_3r+p_4s-c$.}
\label{tab:mode-two}
\begin{tabular}{cc}
\toprule
input wall & mode-two upper closure \\
\midrule
$(2,3,6,5)-7$ & $(2,3,8,5)-7$ \\
$(2,3,6,4)-3$ & $(2,4,6,4)-3$ \\
$(2,3,6,3)-2$ & $(2,4,6,3)-2$ \\
$(2,3,0,2)$   & $(2,3,0,2)$ \\
$(2,5,6,3)-4$ & $(3,5,6,3)-4$ \\
$(2,5,7,4)-5$ & $(3,5,7,4)-5$ \\
$(2,5,9,5)-9$ & $(3,7,9,5)-9$ \\
\bottomrule
\end{tabular}
\end{table}

Let $\overline H_2$ be the maximum of the seven right-hand entries.  Then
\[
   H_2(\alpha)\leq\overline H_2(\alpha).
\]
Applying the same endpoint optimization to mode 3 leaves those seven upper
pieces and adds one new piece,
\[
   4\ell+7m+9r+5s-9.
\]
Thus the joint exact hazard $H_{23}$ obeys
\begin{equation}
\label{eq:joint-hazard}
 H_{23}(\alpha)\leq
 \overline H_{23}(\alpha)
 :=\max\left\{
 \overline H_2(\alpha),
 4\ell+7m+9r+5s-9
 \right\}.
\end{equation}

Only the upper bound is needed for the proof.  Direct $D_4$ character
products attain it throughout the regression simplex, including examples in
which mode 3 genuinely raises the mode-two hazard.  These comparisons protect
against convention or chirality errors, but they are not being extrapolated
to establish the theorem.

\section{The finite low-mode certificate}

\subsection{Five pivot classes}

Denote the $n$th summand of \eqref{eq:telescope} by
\[
 T_n(q)=(1-q)q^{n-1}B(q^n)
 \frac{\prod_{m=1}^{n-1}E(q^m)}
      {\prod_{m=2}^{n}D(q^m)}.
\]
After applying the safe-tail estimates, only five cases remain.

\begin{longtable}{p{0.12\textwidth}p{0.18\textwidth}p{0.27\textwidth}p{0.27\textwidth}}
\caption{Finite data retained in each pivot class.}
\label{tab:pivots}\\
\toprule
pivot & base degree & fermionic core & remaining vector/core factor \\
\midrule
\endfirsthead
\toprule
pivot & base degree & fermionic core & remaining vector/core factor \\
\midrule
\endhead
1 & 0 & none & $(1-q)B(q)$ \\
2 & 1 & $E(q)$ & $B(q^2)/D(q^2)$ \\
3 & 2 & modes 1--2 & vector modes 2--3 \\
4 & 3 & modes 1--3 & vector modes 2--3 \\
$n\geq5$ & at least 4 & modes 1--4 & vector modes 2--3 \\
\bottomrule
\end{longtable}

For $n\geq5$, replacing the actual base degree $n-1$ by 4 is the worst-case
choice.  All omitted fermion modes are at least 5, all omitted vector modes
are at least 4, and $B(q^n)$ is safe.

\subsection{Why first occurrence is sufficient}

Expand each finite fermion core exactly in irreducible $D_4$ characters.  If
the same irrep occurs at several degrees, retain only its first occurrence.
The representation has the same low-vector hazard at every occurrence, while
later copies have more degree available.  This turns the three nontrivial
fermion cores into only
\[
   20,\qquad 50,\qquad 105
\]
distinct first-occurrence records for pivots 3, 4 and $n\geq5$.

For each source $\alpha$ of first degree $f$, the required calculation is
simply
\[
   \text{margin}=(\text{base degree})+f-\overline H_{23}(\alpha).
\]
Pivots 1 and 2 retain the signed coefficients of $B$ explicitly.  For pivot
2 the mode-two upper hazard $\overline H_2$ is used.  The factor $(1-q)$
creates a second copy one degree later, so it can never worsen support.

\subsection{Certificate result}

\begin{table}[ht]
\centering
\caption{Complete low-mode certificate.}
\label{tab:certificate}
\begin{tabular}{crrr}
\toprule
pivot class & records & minimum margin & failures \\
\midrule
1 & 8 & 0 & 0 \\
2 & 140 & 0 & 0 \\
3 & 20 & 0 & 0 \\
4 & 50 & 0 & 0 \\
$n\geq5$ & 105 & 1 & 0 \\
\midrule
total & 323 & 0 & 0 \\
\bottomrule
\end{tabular}
\end{table}

All decompositions are exact Weyl-character calculations over the integers.
The virtual signs in $B$ are retained.  No protected $\Spin(10)$ lift, no
level-1--10 answer and no numerical optimization solver enters the
certificate.

\section{Proof of the all-orders support theorem}

We can now assemble the argument.

\begin{theorem}[All-orders HFT remainder support]
For every $d\geq0$ and every dominant $B_4$ label $\mu$,
\[
   [V_\mu]F_d\neq0
   \quad\Longrightarrow\quad
   d\geq\Kseven(\mu)\geq K(\mu).
\]
\end{theorem}

\begin{proof}
Use the exact telescoping expansion \eqref{eq:telescope} and fix one pivot
$n$.

First decompose the finite fermion core retained in
Table~\ref{tab:pivots}.  For each resulting $D_4$ source, the certificate in
Table~\ref{tab:certificate}, together with
\eqref{eq:joint-hazard}, proves degree at least $\Kseven$ after applying
vector modes 2 and 3.  Pivots 1 and 2 are included directly in the same
certificate.

Every remaining vector mode has index at least 4, so its oscillator degree
pays for the maximum filtered increase.  Every remaining fermion mode has
index at least 5 and is safe for the same reason.  The factor $B(q^n)$ is
safe for $n\geq3$, while its exceptional cases $n=1,2$ were included in the
finite certificate.  Multiplication by $(1-q)$ either preserves the degree or
raises it by one.

Thus every $D_4$ irrep occurring in every coefficient of every pivot summand
has degree at least the folded value of $\Kseven$.  The Weyl-numerator lemma
shows that passage back to dominant $B_4$ characters cannot increase
$\Kseven$.  Hence each individual telescoping summand has the desired
support, and so does their signed sum $F(q)$.

Finally, $\Kseven\geq K$, proving the original Pieri-hazard support theorem.
\end{proof}

\begin{corollary}[Stable-range premise]
Every contribution of the virtual HFT remainder lies in the stable range
required by the consecutive harmonic and symmetric-power Pieri translation
theorem.  The new top-degree contribution also lies outside the translated
exterior-chamber target.
\end{corollary}

The corollary is valid contribution by contribution, including spinorial
sources and negative coefficients in the virtual remainder.

\section{Why the earlier Koszul route failed}

There is a formally correct $B_4$ identity
\[
   F_d=\sum_{j=0}^{9}(-1)^j\Lambda^j(V_9)\otimes R_{d-j}.
\]
It is tempting to bound every tensor product on the right separately using
minimum physical levels.  That strategy is too strong because the sum is
virtual.  An unsafe irrep can occur in one Koszul term and be cancelled by an
equal irrep in another term.

The simplest symptom is already visible near the vacuum: a term
$-V_9q$ violates the desired bound term by term, but is cancelled by the
corresponding positive contribution.  In the standard finite regression,
the generic Koszul-transition test produces 165 such obstructions.  These
are not counterexamples to the theorem; they are counterexamples to a proof
method which discarded essential cancellations.

The Green--Schwarz telescope succeeds because the cancellation required for
division by the universal multiplet is performed algebraically first.  Only
then is a support estimate applied.  This distinction is the conceptual
reason the new proof closes while the principal-embedding shortcut does not.

\section{Independent checks}

The proof is level-independent, but several exact finite comparisons guard
against implementation errors.

\begin{table}[ht]
\centering
\caption{Standard regression checks.  These checks support the implementation;
they are not substitutes for the all-orders argument.}
\label{tab:checks}
\begin{tabularx}{\textwidth}{Xrr}
\toprule
check & cases or terms & failures \\
\midrule
one-mode triality coefficients & 9 & 0 \\
telescoping reconstruction, degrees 0--5 & 6 & 0 \\
separate telescoping support terms & 33 & 0 \\
finite Weyl/Dirac fold regression & 30 & 0 \\
direct low-vector hazard comparisons & 15 & 0 \\
finite low-mode certificate & 323 & 0 \\
observed HFT remainder irreps through physical level 13 & 176 & 0 \\
\bottomrule
\end{tabularx}
\end{table}

The telescope reconstructs the independently derived $B_4$ remainder
weight-for-weight through degree 5.  A larger direct test of the total
$D_4$ irrep support has also passed through degree 24.  These observations
were useful in finding the proof, but the finite-mode reduction is what
removes dependence on a level cutoff.

\section{What has and has not been proved}

\subsection{What is now closed}

The theorem supplies an analytic, all-degree answer to the question which
motivated this report:
\begin{itemize}
  \item higher oscillator levels cannot introduce a new support violation;
  \item chirality folding at the $D_4$ stage cannot increase the required
        $B_4$ degree;
  \item every virtual remainder contribution lies in the stable Pieri range;
  \item the result is baseline-free and does not depend on a selected
        $\Spin(10)$ representative.
\end{itemize}

\subsection{What remains for Bars III.B}

The Bars claim concerns positive integral $\Spin(10)$ lifts of the cumulative
targets $R_n^9$.  The support theorem controls one systematic family of local
inequalities, but the complete finite-horizon lift cone has moving weighted
circuits.  The remaining work is to do at least one of the following:
\begin{enumerate}
  \item prove that the stable Pieri transport assembles all moving
        boundary-tail circuits into a recursively controlled family;
  \item construct an all-level nonnegative integral packetization directly;
  \item prove that another graded construction, such as the BBMS candidate,
        telescopes to a positive cumulative lift at every level.
\end{enumerate}
Even if all rational circuit inequalities are controlled, integrality or
saturation remains a logically separate issue.

\subsection{What the proof says about chirality}

The argument deliberately folds the two $D_4$ chiralities only for the
purpose of a support estimate.  It shows that the support theorem is
insensitive to that choice.  It does \emph{not} prove that different
$\Spin(10)$ chiral lifts are physically equivalent, nor does it select among
the distinct positive completions seen from level 6 onward.  That question
belongs to the later global lift problem.

\section{Reproduction and implementation boundary}

The proof implementation is
\begin{center}
\texttt{trajectory\_kernel/hft\_analytic\_support.py}.
\end{center}
From the project root, the complete standard run is
\begin{verbatim}
./bin/sage-run -m trajectory_kernel.hft_analytic_support \
  --max-level 13 \
  --maximum-label-sum 2 \
  --telescoping-degree 5 \
  --core-pivot 3 --core-pivot 4 --core-pivot 5 --progress
\end{verbatim}
It reports, among other checks,
\begin{verbatim}
Finite low-mode pivots/sources/failures: (3, 4, 5)/175/0
Finite core minimum margin: 0
Complete finite certificate: True
Protected SO(10) baseline imported: NO
\end{verbatim}

The parameter \texttt{--max-level} affects only comparison with explicitly
expanded HFT coefficients.  It does not change the finite proof domain.  The
proof modules do not import any of the protected level-1--10 $\Spin(10)$
baseline files.

\appendix

\section{A compact logical dependency map}

For reference, the proof can be read as the following chain.
\begin{enumerate}
  \item The seven-wall degree $\Kseven$ implies the exact Pieri threshold
        $K$.
  \item Triality gives $E-D=uHB$.
  \item That identity telescopes the Green--Schwarz oscillator product and
        gives equation \eqref{eq:telescope}.
  \item The $D_4$ Weyl numerator maps to $B_4$ without increasing
        $\Kseven$.
  \item Vector and spinor Lipschitz constants close all modes except
        fermions 1--4 and vectors 2--3.
  \item Pieri-polytope optimization bounds vector modes 2--3 by eight affine
        walls.
  \item The remaining fermionic first occurrences form a 323-record exact
        certificate.
  \item Therefore every coefficient of every pivot summand, and hence every
        coefficient of $F(q)$, satisfies degree at least $\Kseven$.
\end{enumerate}

\section{Notation glossary}

\begin{longtable}{p{0.22\textwidth}p{0.70\textwidth}}
\toprule
symbol & meaning \\
\midrule
$B_4$ & $\Spin(9)$, the massive little group in ten dimensions \\
$D_4$ & $\Spin(8)$, used for the transverse Green--Schwarz oscillator product \\
$D_5$ & $\Spin(10)$, the desired parent group for Bars lifts \\
$r_n^9$ & chiral massive-primary $B_4$ character at physical level $n$ \\
$R_n^9$ & cumulative Bars target $\sum_{i=1}^{n}r_i^9$ \\
$R(q)$ & HFT massive-primary series after universal massive-multiplet division \\
$F(q)$ & virtual remainder after removing the level-one $B_4$ vector tower \\
$P_\lambda$ & branching packet $\Res^{D_5}_{B_4}V_\lambda$ \\
$U$ & unitriangular matrix of all-chiral reference packets \\
$c=U^{-1}m$ & integral, generally signed, triangular section coordinates \\
$J$ & global restriction-kernel generator $16_+-16_-$ \\
facet & codimension-one face, or primitive defining inequality, of the packet cone \\
Hilbert basis & minimal integral generators of all lattice points in a rational cone \\
Gale vector & row of an integral kernel-basis matrix for the packet configuration \\
positive circuit & minimally supported nonnegative relation among Gale vectors \\
shell & six-term local kernel identity from consecutive interlacing slices \\
Regge shift & simultaneous $n\mapsto n+1$ and $\ell\mapsto\ell+1$ \\
$\rho=n-\ell$ & elementary invariant of the Regge shift \\
$\delta=n-\ell-3m$ & support-adapted reduced depth; subtracting $3m$ absorbs the level cost of the second Dynkin label \\
upper $\ell$-chamber & high-$\ell$ supported region where multiplicity is inherited from a lower-level Regge predecessor \\
transport boundary & lattice wall separating the upper $\ell$-chamber from the sector where that transport rule is unavailable \\
branching boundary & interlacing face $\mu_1=\lambda_2$ left behind when a branching packet is shifted \\
residual core & supported labels reached by none of the admissible $\ell$-, $m$- or $r$-transport identities \\
stable Pieri range & region $\lambda_1\geq\mu_1$ where one-row tensor multiplicities translate exactly \\
$V,S,C$ & $D_4$ vector, chosen oscillator spinor and opposite spinor \\
$D(u)$ & alternating exterior polynomial of $V$ \\
$E(u)$ & ordinary exterior polynomial of $S$ \\
$B(u)$ & finite triality quotient in $E-D=u(V+S)B$ \\
$K(\mu)$ & exact unstable-Pieri source threshold from TK6B.5e \\
$\Kseven(\mu)$ & uniform seven-wall strengthening of $K$ used in this proof \\
$H_k$ & harmonic one-row representation of degree $k$ \\
pivot $n$ & first oscillator mode selected by the telescoping identity \\
\bottomrule
\end{longtable}

\begin{thebibliography}{9}

\bibitem{Bars}
I.~Bars,
\emph{Stringy evidence for D=11 structure in strongly coupled type-IIA
superstring},
Phys. Rev. D \textbf{52} (1995) 3567--3575,
\href{https://arxiv.org/abs/hep-th/9503228}{arXiv:hep-th/9503228}.

\bibitem{HFT}
A.~Hanany, D.~Forcella and J.~Troost,
\emph{The covariant perturbative string spectrum},
Nucl. Phys. B \textbf{846} (2011) 212--225,
\href{https://arxiv.org/abs/1007.2622}{arXiv:1007.2622}.

\bibitem{GMSZ}
A.~Georgoudis, J.~A.~Minahan, G.~Str\"om and A.~Zoumis,
\emph{The growth of $SO(9)$ super-representations in Type II string theory},
\href{https://arxiv.org/abs/2609.04107}{arXiv:2609.04107}.

\bibitem{Okada}
S.~Okada,
\emph{Pieri rules for classical groups and equinumeration between generalized
oscillating tableaux and semistandard tableaux},
\href{https://arxiv.org/abs/1606.02375}{arXiv:1606.02375}.

\bibitem{BT}
A.~Bourget and J.~Troost,
\emph{The Covariant Chiral Ring},
JHEP \textbf{03} (2016) 163,
\href{https://arxiv.org/abs/1512.03649}{arXiv:1512.03649}.

\bibitem{BM}
T.~Basile and M.~Markou,
\emph{On the deep superstring spectrum},
\href{https://arxiv.org/abs/2405.18467}{arXiv:2405.18467}.

\bibitem{BBMS}
N.~Beisert, M.~Bianchi, J.~F.~Morales and H.~Samtleben,
\emph{On the spectrum of AdS/CFT beyond supergravity},
JHEP \textbf{02} (2004) 001,
\href{https://arxiv.org/abs/hep-th/0310292}{arXiv:hep-th/0310292}.

\end{thebibliography}

\end{document}
